{\centering
	\chapter{Infinite Go}\label{infgo}}
	
\section{Getting used to the board}
\thispagestyle{empty}
\noindent
{\large\textbf{This chapter is a draft, at least for now. I'll include my thought process and most of what I'm able to understand/formalize about it.}}
\vspace{20pt}

To begin with, let's see \textit{where} we are playing now. Our infinite board should be able to have also infinite dimensions, and every one of those could or not be infinite as well.

We'll begin with the simplest case: let's play on a two-dimensional board that has one infinite dimension.

Let $\Gamma$ be a Go game such that,
$$G=\{1,2,\dots,n\}\times \mathbb{Z},$$
for $n\in \mathbb{N}$. We'll use $\mathbb{Z}$ for simplicity as it's a better (easier) representation of the infinite board.
As the number of players does not matter at all, let's use only two players and the canonical adjacency.
Our game will be $\Gamma=(G,E,A_{ort},\Delta_{\mathrm{Go}_{\infty}})$.
At least for now, we are not sure how $\Delta_{\mathrm{Go}_{\infty}}$ works, but we will get there.

Our board could (and for us, will) look like this, Figure~\ref{first_inf_board}. Now, there is a substantial change in this new game.

\begin{figure}[H]
	%\begin{figure}[H]
	\centering
	\begin{tikzpicture}[scale=0.5]
		% Draw the background
		\fill[goban] (-7,-1.5) rectangle (7,6.5);
		% Draw the grid lines
		\foreach \x in {6,5,...,0}{
			\draw[black] (-5,\x) -- (5,\x); 
			% draw horizontal lines
		}
		\foreach \x in {4,...,-4} {
			\draw[black] (\x,0) -- (\x,6); 
			% draw vertical lines
		}
		\foreach \x in {6,0}{
		\draw[very thick] (-5,\x) rectangle (5,\x);
		}% draw upper and lower lines.
		%draw the triple dot line for each infinite side.
		\foreach \x in {6,5,...,0} {
			\node at (-6, \x) {\dots};
			\node at (6, \x) {\dots};
		}
		\foreach \x in {5,...,-5} {
			\node at (\x,-1) {\x};
		}
		\begin{scope}[xshift=15cm,yshift=-0.5cm]
			% Draw the background
			\fill[goban] (-7,-0.5) rectangle (7,6.5);
			% Draw the grid lines
			\foreach \x in {6,5,...,0}{
				\draw[black] (-5,\x) -- (5,\x); 
				% draw horizontal lines
			}
			\foreach \x in {4,...,-4} {
				\draw[black] (\x,0) -- (\x,6); 
				% draw vertical lines
			}
			\foreach \x in {6,0}{
				\draw[very thick] (-5,\x) rectangle (5,\x);
			}% draw upper and lower lines.
			%draw the triple dot line for each infinite side.
			\foreach \x in {6,5,...,0} {
				\node at (-6, \x) {\dots};
				\node at (6, \x) {\dots};
			}
			
			
			\foreach \x in {6,...,0} {
				\filldraw[draw=black,fill=white] (0,\x) circle (12pt);
			}
			
			%draw black stones.
			\foreach \x in {5,...,-5} {
				\filldraw[black] (\x,3) circle (12pt);
			}
			
		\end{scope}
	\end{tikzpicture}
	\caption{A $7\times \infty$ board and two intersecting plays that resolve by definition.}
	\label{first_inf_board}
\end{figure}
Notice that, if the white line were infinite, such play generated two infinite lines in one turn, which is crazy. 


The very first problem on the infinite board is how to determine a capture. Any other prior definition should work just fine.

Let's say we are in the next position, Figure~\ref{first_inf_cap}.
\begin{figure}[H]
	%\begin{figure}[H]
	\centering
	\begin{tikzpicture}[scale=0.5]
		% Draw the background
		\fill[goban] (-7,-0.5) rectangle (7,6.5);
		% Draw the grid lines
		\foreach \x in {6,5,...,0}{
			\draw[black] (-5,\x) -- (5,\x); 
			% draw horizontal lines
		}
		\foreach \x in {4,...,-4} {
			\draw[black] (\x,0) -- (\x,6); 
			% draw vertical lines
		}
		\foreach \x in {6,0}{
			\draw[very thick] (-5,\x) rectangle (5,\x);
		}% draw upper and lower lines.
		%draw the triple dot line for each infinite side.
		\foreach \x in {6,5,...,0} {
			\node at (-6, \x) {\dots};
			\node at (6, \x) {\dots};
		}
		
		%draw black stones.
		\foreach \x in {5,...,-5} {
			\filldraw[black] (\x,3) circle (12pt);
		}
		\begin{scope}[xshift=15cm]
			% Draw the background
			\fill[goban] (-7,-0.5) rectangle (7,6.5);
			% Draw the grid lines
			\foreach \x in {6,5,...,0}{
				\draw[black] (-5,\x) -- (5,\x); 
				% draw horizontal lines
			}
			\foreach \x in {4,...,-4} {
				\draw[black] (\x,0) -- (\x,6); 
				% draw vertical lines
			}
			\foreach \x in {6,0}{
				\draw[very thick] (-5,\x) rectangle (5,\x);
			}% draw upper and lower lines.
			%draw the triple dot line for each infinite side.
			\foreach \x in {6,5,...,0} {
				\node at (-6, \x) {\dots};
				\node at (6, \x) {\dots};
			}
			
			%draw black stones.
			\foreach \x in {5,...,-5} {
				\filldraw[draw=black,fill=white] (\x,4) circle (12pt);
			}
			%draw black stones.
			\foreach \x in {5,...,-5} {
				\filldraw[black] (\x,3) circle (12pt);
			}
			%draw black stones.
			\foreach \x in {5,...,-5} {
				\filldraw[draw=black,fill=white] (\x,2) circle (12pt);
				
			}
		\end{scope}
		
	\end{tikzpicture}
	\caption{An infinite board with an infinite black group being captured.}
	\label{first_inf_cap}
\end{figure}
As this group is infinite, it will touch the very end of the board so we can capture it by surrounding it at the top and bottom.

We can be sure that this constitutes a capture because, denoting the infinite black group by $I$, we have $K(I) \subseteq \bigcup_{W \in S} W$, where $S$ is the set of the two surrounding white groups; that is, $L(I) = \varnothing$.
Now let's take a look at another example, an infinite group but only on one side.
\begin{figure}[H]
	%\begin{figure}[H]
	\centering
	\begin{tikzpicture}[scale=0.5]
		% Draw the background
		\fill[goban] (-7,-0.5) rectangle (7,6.5);
		% Draw the grid lines
		\foreach \x in {6,5,...,0}{
			\draw[black] (-5,\x) -- (5,\x); 
			% draw horizontal lines
		}
		\foreach \x in {4,...,-4} {
			\draw[black] (\x,0) -- (\x,6); 
			% draw vertical lines
		}
		\foreach \x in {6,0}{
			\draw[very thick] (-5,\x) rectangle (5,\x);
		}% draw upper and lower lines.
		%draw the triple dot line for each infinite side.
		\foreach \x in {6,5,...,0} {
			\node at (-6, \x) {\dots};
			\node at (6, \x) {\dots};
		}
		
		%draw black stones.
		\foreach \x in {5,...,0} {
			\filldraw[black] (\x,3) circle (12pt);
		}
		\begin{scope}[xshift=15cm]
			% Draw the background
			\fill[goban] (-7,-0.5) rectangle (7,6.5);
			% Draw the grid lines
			\foreach \x in {6,5,...,0}{
				\draw[black] (-5,\x) -- (5,\x); 
				% draw horizontal lines
			}
			\foreach \x in {4,...,-4} {
				\draw[black] (\x,0) -- (\x,6); 
				% draw vertical lines
			}
			\foreach \x in {6,0}{
				\draw[very thick] (-5,\x) rectangle (5,\x);
			}% draw upper and lower lines.
			%draw the triple dot line for each infinite side.
			\foreach \x in {6,5,...,0} {
				\node at (-6, \x) {\dots};
				\node at (6, \x) {\dots};
			}
			
			%draw white stones.
			\foreach \x in {5,...,0} {
				\filldraw[draw=black,fill=white] (\x,4) circle (12pt);
			}
			%draw black stones.
			\foreach \x in {5,...,0} {
				\filldraw[black] (\x,3) circle (12pt);
			}
			%draw white stones.
			\foreach \x in {5,...,0} {
				\filldraw[draw=black,fill=white] (\x,2) circle (12pt);
				
			}
			\filldraw[draw=black,fill=white] (-1,3) circle (12pt);
		\end{scope}
		
	\end{tikzpicture}
	\caption{An infinite, on one side, black group being captured.}
	\label{first_1sideinf_cap}
\end{figure}
In this case, not using the whole board, even though it's infinite, is better, as it provides one more liberty (or does not lose it).

The following example reflects an interesting nature of infinite Go. (More on this later)

\begin{figure}[H]
	%\begin{figure}[H]
	\centering
	\begin{tikzpicture}[scale=0.5]
		% Draw the background
		\fill[goban] (-7,-0.5) rectangle (7,6.5);
		% Draw the grid lines
		\foreach \x in {6,5,...,0}{
			\draw[black] (-5,\x) -- (5,\x); 
			% draw horizontal lines
		}
		\foreach \x in {4,...,-4} {
			\draw[black] (\x,0) -- (\x,6); 
			% draw vertical lines
		}
		\foreach \x in {6,0}{
			\draw[very thick] (-5,\x) rectangle (5,\x);
		}% draw upper and lower lines.
		%draw the triple dot line for each infinite side.
		\foreach \x in {6,5,...,0} {
			\node at (-6, \x) {\dots};
			\node at (6, \x) {\dots};
		}
		
		%draw black stones.
		\foreach \x in {1,...,-1} {
			\filldraw[black] (\x,3) circle (12pt);
		}
		\foreach \x in {2,4} {
			\filldraw[black] (0,\x) circle (12pt);
		}
		\filldraw[draw=black,fill=white] (0,1) circle (12pt);
		\filldraw[draw=black,fill=white] (0,5) circle (12pt);
		\filldraw[draw=black,fill=white] (-1,4) circle (12pt);
		\filldraw[draw=black,fill=white] (-1,2) circle (12pt);
		\filldraw[draw=black,fill=white] (1,4) circle (12pt);
		\filldraw[draw=black,fill=white] (1,2) circle (12pt);
		\filldraw[draw=black,fill=white] (2,3) circle (12pt);
		\filldraw[draw=black,fill=white] (-2,3) circle (12pt);

		
		\begin{scope}[xshift=15cm]
			% Draw the background
			\fill[goban] (-7,-0.5) rectangle (7,6.5);
			% Draw the grid lines
			\foreach \x in {6,5,...,0}{
				\draw[black] (-5,\x) -- (5,\x); 
				% draw horizontal lines
			}
			\foreach \x in {4,...,-4} {
				\draw[black] (\x,0) -- (\x,6); 
				% draw vertical lines
			}
			\foreach \x in {6,0}{
				\draw[very thick] (-5,\x) rectangle (5,\x);
			}% draw upper and lower lines.
			%draw the triple dot line for each infinite side.
			\foreach \x in {6,5,...,0} {
				\node at (-6, \x) {\dots};
				\node at (6, \x) {\dots};
			}
			
			
			%draw white stones.
			\foreach \x in {5,...,-5} {
				\filldraw[draw=black,fill=white] (\x,4) circle (12pt);
			}
			%draw white stones.
			\foreach \x in {5,...,-5} {
				\filldraw[draw=black,fill=white] (\x,2) circle (12pt);
				
			}
			%draw white stones.
			\foreach \x in {0,1,2,4,5,6} {
				\filldraw[draw=black,fill=white] (-1,\x) circle (12pt);
				
			}
			\foreach \x in {0,1,2,4,5,6} {
				\filldraw[draw=black,fill=white] (1,\x) circle (12pt);
				
			}
			
			%draw black stones.
			\foreach \x in {5,...,-5} {
				\filldraw[black] (\x,3) circle (12pt);
			}
			\foreach \x in {0,...,6} {
				\filldraw[black] (0,\x) circle (12pt);
			}
		\end{scope}
		
	\end{tikzpicture}
	\caption{Two captures.}
	\label{finite_infinite_capture}
\end{figure}
It seems that finite groups are stronger than infinite ones as, again, they don't lose liberties due to not being infinite.

Let's characterize \textit{what} can be done in the game.

\begin{definition}[A move]
	In an infinite game of Go of dimension $n$, the active player of color $e \in E \setminus \{0\}$ at turn $t$ can perform one of the following actions:

		1. Pass: $f(t) = \varnothing$.
		
		2. Propose a connected hyperplane $H \subseteq E \times G$ of dimension at most
		$n-1$. The hyperplane $H$ is a subset such that,
		for at least one dimension, all its points $P(h)$ have a fixed
		coordinate, and all stones $h \in H$ have the color $C(h) = e$ of the player
		making the move. In particular, if $\dim H=0$, then $H$ is a single stone.
		
		The stones already present on the board may intersect the positions of $H$. Let
		\[
		P_{t-1} = \{P(x) \in G \mid x \in T_{t-1},\; C(x) \neq 0\}
		\]
		be the set of occupied points prior to turn $t$. The stones that are
		actually played are then
		\[
		H - P_{t-1} \coloneqq \{h \in H : P(h) \notin P_{t-1}\}.
		\]
		The set $H - P_{t-1}$ may consist of several groups, as the
		stones already present on the board may disconnect the proposed
		hyperplane.
		
		Captures of opposing groups are resolved first; if any connected group contained in $H - P_{t-1}$ remains in a suicide
		condition ($L = \varnothing$), the proposed move is rejected. Otherwise, the
		stones in $H - P_{t-1}$ are placed on the board, and captures
		are resolved according to the usual rules.

\end{definition}
Let me explain why $n-1$ dimensions.

If we look back at the game of Go, stones are placed on a two-dimensional board. In our formulation, this means that we place a connected $0$-dimensional subset of the board, that is, a single stone. In other words, the dimension of the object being played is $n-2$ relative to the dimension of the board.

If we kept this relation, we would obtain a natural generalization for higher-dimensional boards, but we would be unable to define a move in one-dimensional Go, since this would require
$$ n-2=-1. $$
 Thus, in an $n$-dimensional game, a player may propose a connected hyperplane $H$ satisfying
$$ 0\leq\dim H\leq n-1. $$

\section{Infinite SuperKo}

We cannot use Zobrist's method for infinite Go as we would need infinitely many bits. We need a better method; let's try to make one.

If we remember the nature of the SuperKo rule, it says that we cannot replicate a prior board via a stone play. Again, we are using this version of SuperKo (Positional).

Because of the rule, every game has to be a line; that is, we cannot go back into the history of plays (unless it's a pass, which we will not take into consideration here). From the empty board to the very end, we have a straightforward path. Not only that, but because of the rule, the whole game of Go is a \textit{tree}. That is, starting from the empty board, there are 362 possible plays (361 stone placements and a pass) and we cannot go back. From those, every possible outcome, that is, any end of a match, can be found by following a line of unrepeated boards.

Let's define a way to measure how close two given boards are.

First, let's notice that, in the very beginning, there is exactly one region and not much else on the board. That is, before the first turn, $|I^0|=1$ (the big empty region).

Now, the very first turn will make $|I^0|=1$ if the black player plays a stone or a finite line. If the player chooses to play an infinite one, then $|I^0|=2$. Let's pretend they just played an infinite group; with that, Zobrist's method or any method that uses the stone itself goes nuts and can't be used here, but $|I^0|$ is just a little bigger, finitely bigger. Now, this is not perfect as there are infinitely many ways to play an infinite group so we end up with $|I^0|=2$. Before going forward, let's see an example.

\begin{figure}[H]
	%\begin{figure}[H]
	\centering
	\begin{tikzpicture}[scale=0.5]
		% Draw the background
		\fill[goban] (-7,-0.5) rectangle (7,6.5);
		\filldraw[fill=orange, fill opacity=0.35, draw=black] (-5.35,0) rectangle (5.35,6);
		% Draw the grid lines
		\foreach \x in {6,5,...,0}{
			\draw[black] (-5,\x) -- (5,\x); 
			% draw horizontal lines
		}
		\foreach \x in {4,...,-4} {
			\draw[black] (\x,0) -- (\x,6); 
			% draw vertical lines
		}
		\foreach \x in {6,0}{
			\draw[very thick] (-5,\x) rectangle (5,\x);
		}% draw upper and lower lines.
		%draw the triple dot line for each infinite side.
		\foreach \x in {6,5,...,0} {
			\node at (-6, \x) {\dots};
			\node at (6, \x) {\dots};
		}
		
		%draw black stones.
%		\foreach \x in {5,...,-5} {
%			\filldraw[black] (\x,3) circle (12pt);
%		}
		\begin{scope}[xshift=15cm]
			% Draw the background
			\fill[goban] (-7,-0.5) rectangle (7,6.5);
			% Draw the grid lines
			\filldraw[fill=violet, fill opacity=0.35, draw=black] (-5.35,3.65) rectangle (5.35,6);
			\filldraw[fill=violet, fill opacity=0.35, draw=black] (-5.35,0) rectangle (5.35,2.35);
			
			
			\foreach \x in {6,5,...,0}{
				\draw[black] (-5,\x) -- (5,\x); 
				% draw horizontal lines
			}
			\foreach \x in {4,...,-4} {
				\draw[black] (\x,0) -- (\x,6); 
				% draw vertical lines
			}
			\foreach \x in {6,0}{
				\draw[very thick] (-5,\x) rectangle (5,\x);
			}% draw upper and lower lines.
			%draw the triple dot line for each infinite side.
			\foreach \x in {6,5,...,0} {
				\node at (-6, \x) {\dots};
				\node at (6, \x) {\dots};
			}
			
			
			%draw black stones.
			\foreach \x in {5,...,-5} {
				\filldraw[black] (\x,3) circle (12pt);
			}
			
		\end{scope}
		
	\end{tikzpicture}
	\caption{A board where the number of regions merely grows by one when an infinite group is played.}
	\label{inf_superko_intuition}
\end{figure}
As I said earlier, playing such a group vertically (in this case it will be finite, but you get the idea) or shifting it up or down, will end up with the same number of regions.

The first problem is to determine if the number of regions is always finite or not.

The second problem is how to determine that boards having the same number of regions are in fact different; we cannot look at them point by point as that's too costly. My idea is, if we can solve the first problem, then we can measure every region and, by that, we can tell that even if two boards have the same number of regions, those regions might not be equal (by biggest and smallest radius or something like that). 

The third problem is color shifting: two boards might not be equal, and the regions could be exactly the same (number and size), but the colors of the groups surrounding them might not. (I'm thinking about discriminating between finite and infinite regions. Then, in the set of finite ones, we can calculate the influence of each color; this should work, aunque it's just another filter and is far from perfect.)


 Let's attack the first problem, but we need to be a little more precise about what we want.
 
Given an infinite game of Go, is the number of regions always finite? If there are infinitely many moves, then the answer is no. Let's consider infinite one-dimensional Go.

\begin{figure}[H]
	%\begin{figure}[H]
	\centering
	\begin{tikzpicture}[scale=0.7]
		% Draw the background
		\fill[goban] (-10,-1.5) rectangle (10,0.5);
		\foreach \x in {0} {
			\node at (-9, \x) {\dots};
			\node at (9, \x) {\dots};
		}
		
		\foreach \x in {0,1,...,6}{
			\draw[black] (\x,-0.2) -- (\x,0.2); 
			% draw horizontal lines
		}
		
		\draw[very thick] (-8.5,0) rectangle (8.5,0);
		
		% draw vertical lines
		\foreach \x in {-8,...,8} {
			\draw[black] (\x,-0.25) -- (\x,0.25); 
			
		}
		\foreach \x in {8,...,-8} {
			\node at (\x,-1) {\x};
		}

		\begin{scope}[yshift=-2.5cm]
		% Draw the background
		\fill[goban] (-10,-0.5) rectangle (10,0.5);
		\foreach \x in {0} {
			\node at (-9, \x) {\dots};
			\node at (9, \x) {\dots};
		}
		
		\foreach \x in {0,1,...,6}{
			\draw[black] (\x,-0.2) -- (\x,0.2); 
			% draw horizontal lines
		}
		
		\draw[very thick] (-8.5,0) rectangle (8.5,0);
		
		% draw vertical lines
		\foreach \x in {-8,...,8} {
			\draw[black] (\x,-0.25) -- (\x,0.25); 
			
		}
		\foreach \x in {-8,-6,-4,-2,0,2,4,6,8} {
			\filldraw[fill=violet, fill opacity=0.35, draw=black] (\x-0.3,-0.3) rectangle (\x+0.3,0.3);
		}
		\foreach \x in {-7,-5,-3,-1,1,3,5,7} {
			\filldraw[black] (\x,0) circle (12pt);
		}
		\end{scope}
	\end{tikzpicture}
	\caption{A one-dimensional infinite Go board and an example of infinitely many \textcolor{violet}{regions} produced by infinitely many plays.}
	\label{inf_onedim}
	
	\begin{goban}[19x1, partial=A1:W1, label=false,scale=1.6]
		\pgfgoticks
	\end{goban}
\end{figure}
It is possible to construct many examples like this that end up with infinitely many regions. The trick here—and perhaps our salvation—is that there are infinitely many plays, which, given the turn-by-turn nature of the game, shouldn't happen at any finite stage. If we restrict ourselves to finitely many plays, then our question becomes:

Given an infinite game of Go, is the number of regions generated by a finite number of plays always finite too?
More precisely, we want to prove that:
$$t < \infty \implies |I^0(T_t)|<\infty.$$

In order to do this, we need to show something about the nature of the game; we'll characterize how to split a region. For that, we need some tools first.

\begin{definition}[Dense hyperplane on a region]
	Given a board $T_t$, a region $R \in I^0(T_{t-1})$, and a proposed hyperplane $H$, we say that $H$ is \emph{dense} across $R$ in $T_t$
	if every intersection $p \in P(H) \cap P(R)$ is occupied in $T_t$, that is,
	$$(0, p) \notin T_t \quad \forall p \in P(H) \cap P(R).$$
\end{definition}

\begin{lemma}[Adjacency bound]\label{lem_2n}
	Given an $n$-dimensional board of Go ($n \in \mathbb{N}$) with the canonical orthogonal adjacency where each axis has size at least $2$, a point $t$ on the board has at most $2n$ adjacencies and at least $n$:
	$$n \leq |K(t)| \leq 2n.$$
	If all $n$ dimensions are bi-infinite ($G = \mathbb{Z}^n$), then $|K(t)| = 2n$ for every $t \in G$.
\end{lemma}
\begin{proof}
	This follows directly from the coordinate-wise analysis in Lemma~\ref{lem:2n}.
\end{proof}


%not yet
% \begin{definition}[Interior of a region]
	%	Given a region $R$, its interior $R^\circ$ are the points in $R$ such that its adjacency set its in $R$, that is:
	%	$$R^\circ = \{r \in R: K(r)\subset R \}$$
	% \end{definition}






%this is the proper theorem, just think in a 1 dim region on a two dim board (a line) then a single stone (0 dim) can split it (if in the interior)
\begin{theorem}\label{lem_bigplays}
	In an $n$-dimensional game of Go, a single move played on the empty board $T_0$ that splits an unbounded region $R$ of infinite dimension $d$ into two disconnected components must be a dense hyperplane of dimension at least $d-1$ along the unbounded axes (and in general, any separator splitting a $d$-dimensional region must have cumulative boundary dimension at least $d-1$).
\end{theorem}
This tells us \textit{how} to split a region, and we can now prove our main finiteness result.

%we need this to proove that the number of regions are always finite. (I hope)

\begin{figure}[H]
	%\begin{figure}[H]
	\centering
	\begin{tikzpicture}[scale=0.7]
		\node at (-11,0) {$t=0$};
		\fill[goban] (-10,-0.5) rectangle (10,0.5);
		\foreach \x in {0} {
			\node at (-9, \x) {\dots};
			\node at (9, \x) {\dots};
		}
		
		\foreach \x in {0,1,...,6}{
			\draw[black] (\x,-0.2) -- (\x,0.2); 
			% draw horizontal lines
		}
		
		\draw[very thick] (-8.5,0) rectangle (8.5,0);
		
		% draw vertical lines
		\foreach \x in {-8,...,8} {
			\draw[black] (\x,-0.25) -- (\x,0.25); 
			
		}
		%draw a region
		\filldraw[fill=orange, fill opacity=0.35, draw=black] (-8.65,-0.3) rectangle (8.65,0.3);
		
		\begin{scope}[yshift=-1.5cm]
		\node at (-11,0) {$t=1$};
			% Draw the background
			\fill[goban] (-10,-0.5) rectangle (10,0.5);
			\foreach \x in {0} {
				\node at (-9, \x) {\dots};
				\node at (9, \x) {\dots};
			}
			
			\foreach \x in {0,1,...,6}{
				\draw[black] (\x,-0.2) -- (\x,0.2); 
				% draw horizontal lines
			}
			
			\draw[very thick] (-8.5,0) rectangle (8.5,0);
			
			% draw vertical lines
			\foreach \x in {-8,...,8} {
				\draw[black] (\x,-0.25) -- (\x,0.25); 
				
			}
			%draw a region
			\filldraw[fill=violet, fill opacity=0.35, draw=black] (-8.65,-0.3) rectangle (1.35,0.3);
			\filldraw[fill=violet, fill opacity=0.35, draw=black] (2.65,-0.3) rectangle (8.65,0.3);
			
			
			%draw black stones.
			\foreach \x in {2} {
				\filldraw[black] (\x,0) circle (12pt);
			}
		\end{scope}
	\end{tikzpicture}
	\caption{An infinite one-dimensional board and the first turn showing how a region is split.}
	\label{inf_ind_n=1_t=1}
\end{figure}

\begin{theorem}[Finite regions induced by finite plays]
	Given an $n$-dimensional infinite game of Go starting from the empty board $T_0$, the number of empty regions generated by a finite number of plays $t \in \mathbb{N}$ is always finite:
	$$t < \infty \implies |I^0(T_t)| < \infty.$$
\end{theorem}
\begin{proof}
	We proceed by induction on the number of turns $t \ge 0$ for any fixed dimension $n \ge 1$.
	
	For $n=1$ and $t=1$, starting from the single empty region at $t=0$, we can only play single stones ($\dim H = 0$). Placing a stone splits the region into at most two regions (Figure~\ref{inf_ind_n=1_t=1}).
	
	In general, for any $n \ge 1$:
	\begin{itemize}
		\item \textbf{Base case ($t = 0$):} The board begins empty and connected, so there is exactly one region ($|I^0(T_0)| = 1 < \infty$) and zero stone groups ($|\mathcal{G}(T_0)| = 0$).
		\item \textbf{Inductive step ($t - 1 \to t$):} Assume that at turn $t - 1$, both the number of empty regions $p = |I^0(T_{t-1})|$ and the number of stone groups $|\mathcal{G}(T_{t-1})|$ are finite.
		
		At turn $t$, if the player passes ($f(t) = \varnothing$), then $T_t = T_{t-1}$ and the number of regions remains $p < \infty$. Otherwise, a legal play proposes a connected affine hyperplane $H_t$ of dimension at most $n - 1$. Since $H_t$ is defined by fixing at least one coordinate, intersecting each of the $p$ existing regions in $I^0(T_{t-1})$ with the complement of $H_t$ splits each region into at most a finite number of connected components (the two half-spaces on either side of $H_t$, plus at most finitely many bounded cavities enclosed together with the finitely many prior hyperplanes $H_1, \dots, H_{t-1}$).
		
		Furthermore, resolving captures can at most remove a subset of the finitely many existing groups in $\mathcal{G}(T_{t-1})$, turning their stones into vacant intersections that either form finitely many new regions or merge adjacent existing regions. Therefore, at turn $t$, the number of regions $|I^0(T_t)|$ and stone groups $|\mathcal{G}(T_t)|$ remains finite.
	\end{itemize}
\end{proof}

%this might not be usefull at all.
%\begin{definition}[Color difference function]
%	Given two boards $T_1$ and $T_2$ of the same Go game, the difference at a point $t=(t_1,...t_m)$ it's define as,
%	\[ 
%	\delta(T_1,T_2)_t=
%	\begin{cases}
%		1 \quad \text{ if }C(t)_{T_1} \neq C(t)_{T_2}, \\
%		0 \quad else. \\
%	\end{cases}
%	\]
%\end{definition}
%Now, let's apply this to every point.
%\begin{definition}[Distance between Tables]
%	Given two boards $T_1$ and $T_2$ of the same Go game, the difference between them will be:
%	$$\sum_{t\in G} \delta(T_1,T_2)_{t}.$$
%\end{definition}
%
%The method should be simple, if the distance is $0$ that mean they are exactly the same board. Now, let's say we try to compute this, we might found that for the very first millionth positions they are equal but, we have infinite of those.
%I have no idea how to use this, yet.
